\begin{exercise}{6.1}{5}
    Man betrachte in der komplexen Zahlenebene die Kreisspiegelung am
    Kreis mit Radius $2$ und Zentrum in $i$. Was ist das Bild von $2$? Was ist das Bild des
    Einheitskreises?
\end{exercise}

Die Kreisspiegelung am Kreis $K=\{z\in\mathbb C\mid |z-i|=2\}$
ist nach 2.1.3. gegeben durch (für alle $w \in \mathbb C$ gilt $w\overline w=|w^2|$)
\[
    \sigma(z) \quad = \quad i+\frac{4}{|z-i|^2}(z-i) \quad = \quad i+\frac{4}{\overline{z-i}}.
\]
Zunächst berechnen wir das Bild des Punktes $2$:
\begin{align*}
\sigma(2)
&=
i+\frac{4}{\overline{2-i}}\\
&=
i+\frac{4}{2+i}\\
&=
i+\frac{4(2-i)}{(2+i)(2-i)}\\
&=
i+\frac{8-4i}{5}\\
&=\boxed{\frac85+\frac15 i.}
\end{align*}
Nun bestimmen wir das Bild des Einheitskreises
\[
S^1=\{z\in\mathbb C\mid |z|=1\}.
\]
Laut Lemma 2.1.6 muss das Bild von $S^1$ wieder ein verallgemeinerter Kreis sein.
Da das Zentrum $i$ auf dem Einheitskreis liegt, ist $\infty$ im Bild des Einheitskreises und nach Def. 2.1.2 ist das Bild somit eine erweiterte Gerade.
Das Berechnen von Spiegelungen zweier Punkte $\neq i$ des Einheitskreises liefert somit das gesuchte Bild.
Die Berechnung von $\sigma(1)$ und $\sigma(-1)$ ergibt nach analogen Schritten wie oben bei $\sigma(2)$:
\[
\sigma(1)=i+\frac{4}{\overline{1-i}}=2-i
\]
und
\[
\sigma(-1)=i+\frac{4}{\overline{-1-i}}=-2-i
\]
Damit ist das Bild des Einheitskreises die erweiterte Gerade
\[
\boxed{\{\,w\in\mathbb C\mid \operatorname{Im}(w)=-1\,\} 
\sqcup \{\infty\}.}
\]